\section*{§2. Второе начало. Энтропия}
\addcontentsline{toc}{section}{§2. Второе начало. Энтропия}

\begin{basics}
\sub{Энтропия}
\[\dd S=\frac{\delta Q}{T}\quad\text{(только для обратимого процесса)},\qquad S_2-S_1=\int_1^2\frac{\delta Q}{T}.\]
$S$ — функция состояния (Дж/К): $\Delta S$ зависит только от начального и конечного состояний, не от пути.
$1/T$ — интегрирующий множитель: $\delta Q$ — не полный дифференциал, а $\delta Q/T$ — полный.

\sub{Второе начало}
\begin{itemize}
\item Неравенство Клаузиуса: $\displaystyle\oint\frac{\delta Q}{T}\le0$; «$=$» для обратимого цикла, «$<$» для необратимого.
\item В изолированной (теплоизолированной) системе $\Delta S\ge0$: при необратимых процессах (смешение, теплообмен при
разных $T$, трение, расширение в пустоту, резкий скачок давления) энтропия \emph{растёт}, при обратимых — постоянна.
\item КПД любой машины между $T_1$ и $T_2$ не больше, чем у Карно: $\eta\le1-T_2/T_1$.
\end{itemize}

\sub{Как считать $\Delta S$ в необратимом процессе}
Формула $\dd S=\delta Q/T$ применима только к обратимому пути. Поэтому:
(1) найти конечное равновесное состояние (из закона сохранения энергии и условий равновесия);
(2) соединить начало и конец \emph{любым удобным обратимым} путём и посчитать $\Delta S$ по нему.
Для идеального газа это уже сделано в формуле ниже.

\sub{Энтропия идеального газа}
Из $T\dd S=\nu C_V\dd T+p\,\dd V$, $p=\nu RT/V$:
\begin{align*}
\Delta S&=\nu\Big[C_V\ln\frac{T_2}{T_1}+R\ln\frac{V_2}{V_1}\Big]
=\nu\Big[C_p\ln\frac{T_2}{T_1}-R\ln\frac{p_2}{p_1}\Big]\\
&=\nu\Big[C_V\ln\frac{p_2}{p_1}+C_p\ln\frac{V_2}{V_1}\Big].
\end{align*}
Частные случаи: изотерма $\nu R\ln\frac{V_2}{V_1}$, изохора $\nu C_V\ln\frac{T_2}{T_1}$,
изобара $\nu C_p\ln\frac{T_2}{T_1}$, обратимая адиабата $0$, политропа $\nu C\ln\frac{T_2}{T_1}$.
Энтропия аддитивна: у смеси/системы $\Delta S=\sum_i\Delta S_i$, каждая компонента считается так, будто она одна
занимает весь свой объём.

\sub{Конденсированные тела и фазовые переходы}
Нагрев тела с теплоёмкостью $C$ ($=mc$): $\Delta S=\int\frac{C\,\dd T}{T}=C\ln\frac{T_2}{T_1}$.
\\Плавление/кипение при постоянной $T$: $\Delta S=\frac{Q}{T}=\frac{\lambda m}{T}$, $\lambda$ — удельная теплота перехода.

\sub{Диаграмма $T$--$S$}
$Q=\int T\,\dd S$ — площадь под кривой; $A$ за цикл — площадь цикла; изотерма — горизонталь, обратимая адиабата — вертикаль.

\sub{Полезное тождество}
$\Big(\dfrac{\partial U}{\partial V}\Big)_T=T\Big(\dfrac{\partial p}{\partial T}\Big)_V-p$ (из $\dd U=T\dd S-p\,\dd V$); для идеального газа даёт $0$.
\end{basics}

%-------------------------------------------------------------
\begin{problem}{2.1}{ex}{Азот, занимающий объём $V_1=1$\,м$^3$ при давлении $p=10^5$\,Па и температуре $T=300$\,К, расширяется изотермически
до объёма $V_2=2$\,м$^3$. Вычислить изменение энтропии газа $\Delta S$.}
\textbf{1.} Изотерма: $\Delta S=\nu R\ln\frac{V_2}{V_1}$ (формула энтропии ид.\ газа при $T_2=T_1$).

\textbf{2.} Число молей из уравнения состояния: $\nu R=\frac{pV_1}{T}$. Тогда
$\Delta S=\frac{pV_1}{T}\ln\frac{V_2}{V_1}=\frac{10^5\cdot1}{300}\cdot\ln2=333\cdot0{,}693$.
\answer{$\Delta S\approx231$ Дж/К.}
\end{problem}

%-------------------------------------------------------------
\begin{problem}{2.2}{ex}{Найти изменение энтропии льда массой $m=0{,}03$\,кг при превращении его в пар, если начальная температура льда
$-40^\circ$C, а температура пара $100^\circ$C. Удельные теплоёмкости льда и воды $c_1=2{,}1$ и $c_2=4{,}2$\,кДж/(кг\,К);
теплота плавления $\lambda_1=330$\,кДж/кг, парообразования $\lambda_2=2{,}25\cdot10^3$\,кДж/кг. Давление атмосферное.}
\textbf{Идея:} разбить процесс на этапы и сложить изменения энтропии (энтропия аддитивна).
\begin{enumerate}
\item Нагрев льда $233\to273$\,К: $\Delta S_1=mc_1\ln\frac{273}{233}=0{,}03\cdot2100\cdot0{,}158=10$\,Дж/К.
\item Плавление при $273$\,К: $\Delta S_2=\frac{m\lambda_1}{273}=\frac{0{,}03\cdot3{,}3\cdot10^5}{273}=36{,}2$\,Дж/К.
\item Нагрев воды $273\to373$\,К: $\Delta S_3=mc_2\ln\frac{373}{273}=0{,}03\cdot4200\cdot0{,}312=39{,}3$\,Дж/К.
\item Кипение при $373$\,К: $\Delta S_4=\frac{m\lambda_2}{373}=\frac{0{,}03\cdot2{,}25\cdot10^6}{373}=181$\,Дж/К.
\end{enumerate}
\answer{$\Delta S=\Delta S_1+\Delta S_2+\Delta S_3+\Delta S_4\approx266{,}5$ Дж/К.}
\end{problem}

%-------------------------------------------------------------
\begin{problem}{2.3}{exdz}{Определить работу, совершаемую произвольным телом за обратимый цикл из изотермы 1--2, политропы 2--3 и адиабаты 3--1.
Известны теплоёмкость $C$ тела в политропном процессе и температуры $T_1$ (на изотерме) и $T_3$.}
\textbf{Идея:} уравнение состояния тела неизвестно, поэтому теплоту на изотерме находим из второго начала.

\textbf{1. Равенство Клаузиуса} для обратимого цикла $\oint\frac{\delta Q}{T}=0$, по участкам:
\[\underbrace{\frac{Q_{12}}{T_1}}_{\text{изотерма}}+\underbrace{\int_{T_1}^{T_3}\frac{C\,\dd T}{T}}_{\text{политропа}}+\underbrace{0}_{\text{адиабата}}=0
\ \Rightarrow\ Q_{12}=-CT_1\ln\frac{T_3}{T_1}=CT_1\ln\frac{T_1}{T_3}.\]

\textbf{2. Теплота на политропе:} $Q_{23}=C(T_3-T_1)$; на адиабате $Q_{31}=0$.

\textbf{3. Работа за цикл} $A=\sum Q$ (за цикл $\Delta U=0$).
\answer{$A=C\Big[T_1\ln\dfrac{T_1}{T_3}+T_3-T_1\Big]$.}
\end{problem}

%-------------------------------------------------------------
\begin{problem}{2.4}{ex}{Показать, что для любого вещества политропа может пересекать изотерму не более чем в одной точке.}
\textbf{От противного.} Пусть политропа и изотерма $T_0$ пересекаются в точках $A$ и $B$. Составим цикл: $A\to B$ по политропе, $B\to A$ по изотерме.
\begin{enumerate}
\item На политропе ($C=$\,const): $Q_{\text{пол}}=C(T_B-T_A)=0$ и $\int\frac{C\,\dd T}{T}=C\ln\frac{T_B}{T_A}=0$, так как $T_A=T_B=T_0$.
\item Клаузиус для обратимого цикла: $0+\frac{Q_{\text{из}}}{T_0}=0\Rightarrow Q_{\text{из}}=0$.
\item Тогда работа за цикл $A=Q_{\text{пол}}+Q_{\text{из}}=0$. Но кривые разные, цикл охватывает ненулевую площадь на $p$--$V$, т.е.\ $A\ne0$. Противоречие.
\end{enumerate}
\end{problem}

%-------------------------------------------------------------
\begin{problem}{2.5}{ex}{Электромагнитное излучение в полости объёмом $V$ находится в равновесии со стенками при температуре $T$.
Его энергия $U=u(T)V$, давление $p=u/3$. Найти связь $u$ с $T$ (закон Стефана–Больцмана).}
\textbf{1. Тождество} из $\dd U=T\dd S-p\,\dd V$ (следствие второго начала):
$\big(\frac{\partial U}{\partial V}\big)_T=T\big(\frac{\partial p}{\partial T}\big)_V-p$.

\textbf{2. Подставляем:} $\big(\frac{\partial U}{\partial V}\big)_T=u$, $\big(\frac{\partial p}{\partial T}\big)_V=\frac{u'}3$:
\quad $u=\frac T3u'-\frac u3\Rightarrow4u=T\frac{\dd u}{\dd T}$.

\textbf{3. Разделяем переменные:} $\frac{\dd u}{u}=4\frac{\dd T}{T}\Rightarrow\ln u=4\ln T+\text{const}$.
\answer{$u=\sigma T^4$.}
\end{problem}

%-------------------------------------------------------------
\begin{problem}{2.6}{exdz}{В качестве переменных состояния выберем $T$ и $S$. Построить цикл Карно на диаграмме $T$--$S$ и найти с её помощью КПД.}
\textbf{1. Вид цикла.} Изотермы $T_1$ и $T_2$ — горизонтальные прямые. Обратимые адиабаты: $\delta Q=0\Rightarrow\dd S=0$ — вертикальные прямые $S_1$ и $S_2$.
Цикл — прямоугольник.

\textbf{2. Теплоты.} На $T$--$S$ теплота $Q=\int T\,\dd S$ — площадь под кривой:
на верхней изотерме получено $Q_1=T_1(S_2-S_1)$, на нижней отдано $|Q_2|=T_2(S_2-S_1)$.

\textbf{3. КПД:} $\eta=\frac{Q_1-|Q_2|}{Q_1}$, $(S_2-S_1)$ сокращается.
\answer{$\eta=\dfrac{T_1-T_2}{T_1}$ — для любого рабочего тела.}
\end{problem}

%-------------------------------------------------------------
\begin{problem}{2.7}{ex}{Система квазистатически переводится из состояния 1 в состояние 2 двумя способами: (а) адиабатически охлаждается до $T_0$,
изотермически получает тепло, затем адиабатически переходит в 2; (б) по пути, на каждом участке которого система получает тепло, а её температура выше $T_0$.
В каком случае система получает больше тепла?}
\textbf{1.} Изменение энтропии $S_2-S_1$ одинаково (функция состояния), и $S_2>S_1$, так как в (б) всё время $\delta Q>0$.

\textbf{2. Способ (а):} тепло получено только на изотерме: $Q_a=T_0(S_2-S_1)$.

\textbf{3. Способ (б):} $Q_b=\int T\,\dd S$, и везде $T>T_0$, $\dd S>0$, поэтому $Q_b>T_0\int\dd S=T_0(S_2-S_1)$.
\answer{$Q_b>Q_a$.}
\end{problem}

%-------------------------------------------------------------
\begin{problem}{2.8}{ex}{Теплоизолированный сосуд разделён невесомым поршнем пополам. С одной стороны $\nu$ молей идеального газа при $T_0$, $p_0$,
с другой — вакуум. Поршень отпускают, газ заполняет весь сосуд. Затем объём медленно возвращают к исходному. Найти изменения внутренней энергии и энтропии.}
\textbf{1. Расширение в пустоту.} Поршень невесомый — газ работы не совершает ($A=0$), сосуд теплоизолирован ($Q=0$), значит $\Delta U=0$ и $T=T_0$.
Процесс необратимый, но $\Delta S$ считаем по изотерме между теми же состояниями: $\Delta S_1=\nu R\ln2$.

\textbf{2. Медленное сжатие} в теплоизолированном сосуде — обратимая адиабата: $\Delta S_2=0$;
$TV^{\gamma-1}=\text{const}\Rightarrow T=T_0\cdot2^{\gamma-1}$.

\textbf{3. Итог:} $\Delta U=\nu C_V(T-T_0)$, $\Delta S=\Delta S_1+\Delta S_2$.
\answer{$\Delta U=\nu C_VT_0\big(2^{\gamma-1}-1\big)$,\; $\Delta S=\nu R\ln2$.}
\end{problem}

%-------------------------------------------------------------
\begin{problem}{2.9}{exdz}{Приводимые в тепловой контакт два тела имеют одинаковые теплоёмкости $C_p=\text{const}$ и разные температуры $T_1$ и $T_2$.
Найти приращение энтропии в результате установления теплового равновесия.}
\textbf{1. Конечная температура} — из сохранения энергии: $C_p(T-T_1)+C_p(T-T_2)=0\Rightarrow T=\frac{T_1+T_2}2$.

\textbf{2. Энтропия каждого тела} при нагреве/охлаждении: $\Delta S_i=C_p\ln\frac{T}{T_i}$. Складываем:
$\Delta S=C_p\ln\frac{T^2}{T_1T_2}$.
\answer{$\Delta S=C_p\ln\dfrac{(T_1+T_2)^2}{4T_1T_2}$.}
\note{$\Delta S>0$, т.к.\ $(T_1+T_2)^2\ge4T_1T_2$ (среднее арифметическое $\ge$ среднего геометрического) — процесс необратим.}
\end{problem}

%-------------------------------------------------------------
\begin{problem}{2.10}{exdz}{В двух одинаковых сосудах находятся различные идеальные газы: масса первого $M_1$ (молярная масса $\mu_1$), второго — $M_2$ ($\mu_2$).
Температуры одинаковы. Сосуды соединяют, начинается диффузия. Найти суммарное изменение энтропии.}
\textbf{1. Конечное состояние.} Газы идеальные — не взаимодействуют; энергия сохраняется, температура не меняется.
Каждый газ просто расширяется на весь объём: $V\to2V$.

\textbf{2. Энтропия каждого газа} — как при изотермическом расширении: $\Delta S_i=\nu_iR\ln2=\frac{M_i}{\mu_i}R\ln2$.
\answer{$\Delta S=\Big(\dfrac{M_1}{\mu_1}+\dfrac{M_2}{\mu_2}\Big)R\ln2$.}
\note{парадокс Гиббса: для \emph{одинаковых} газов формула дала бы $\Delta S>0$, хотя состояние не изменилось. Для тождественных газов $\Delta S=0$.}
\end{problem}

%-------------------------------------------------------------
\begin{problem}{2.11}{ex}{Какую максимальную работу можно получить, используя два тела с температурами $T_{10}$ и $T_{20}$ ($T_{10}>T_{20}$) как нагреватель и холодильник?
Теплоёмкости $C_1$ и $C_2$ постоянны. Найти конечную температуру $T$.}
\textbf{1.} Работа максимальна, если процесс обратимый (машина Карно с бесконечно малыми циклами); тогда энтропия системы не меняется:
\[C_1\ln\frac{T}{T_{10}}+C_2\ln\frac{T}{T_{20}}=0\ \Rightarrow\ T^{C_1+C_2}=T_{10}^{C_1}T_{20}^{C_2}.\]

\textbf{2. Работа} = убыль энергии тел: $A=C_1(T_{10}-T)+C_2(T_{20}-T)$.
\answer{$T=T_{10}^{C_1/(C_1+C_2)}\,T_{20}^{C_2/(C_1+C_2)}$,\; $A_{\max}=C_1T_{10}+C_2T_{20}-(C_1+C_2)T$.}
\end{problem}

%-------------------------------------------------------------
\begin{problem}{2.16}{dz}{Вычислить КПД цикла, изображённого на рис.~2.6 (диаграмма $T$--$S$). Известны $T_1$, $T_2$, $S_1$, $S_2$ и заштрихованные площади
$\Delta Q_1$ и $\Delta Q_2$.}
На рисунке цикл вписан в прямоугольник $T_2\le T\le T_1$, $S_1\le S\le S_2$; заштрихованные области — части прямоугольника, \emph{не} входящие в цикл:
$\Delta Q_1$ — у верхней стороны, $\Delta Q_2$ — у нижней.

\textbf{1.} На $T$--$S$ теплота равна площади под кривой. Верхняя ветвь (вправо, $\dd S>0$) — тепло получено. Она проходит ниже стороны $T_1$ на площадь $\Delta Q_1$:
$Q_1=T_1(S_2-S_1)-\Delta Q_1$.

\textbf{2.} Нижняя ветвь (влево) — тепло отдано. Она проходит выше стороны $T_2$ на площадь $\Delta Q_2$:
$|Q_2|=T_2(S_2-S_1)+\Delta Q_2$.

\textbf{3.} $A=Q_1-|Q_2|=(T_1-T_2)(S_2-S_1)-\Delta Q_1-\Delta Q_2$ (площадь цикла), $\eta=A/Q_1$.
\answer{$\eta=\dfrac{(T_1-T_2)(S_2-S_1)-\Delta Q_1-\Delta Q_2}{T_1(S_2-S_1)-\Delta Q_1}$.}
\end{problem}

%-------------------------------------------------------------
\begin{problem}{2.19}{dz}{Цикл состоит из изобары 1--2, изохоры 2--3 и адиабаты 3--1. Показать, что для любого вещества с постоянными молярными теплоёмкостями
$C_V$ и $C_p$ температуры связаны соотношением $T_2/T_3=(T_2/T_1)^\gamma$, $\gamma=C_p/C_V$.}
\textbf{1.} Энтропия — функция состояния, за цикл $\Delta S=0$. Изменения по участкам:
изобара $\Delta S_{12}=C_p\ln\frac{T_2}{T_1}$, изохора $\Delta S_{23}=C_V\ln\frac{T_3}{T_2}$, адиабата $\Delta S_{31}=0$.

\textbf{2.} $C_p\ln\frac{T_2}{T_1}+C_V\ln\frac{T_3}{T_2}=0$. Делим на $C_V$: $\gamma\ln\frac{T_2}{T_1}=\ln\frac{T_2}{T_3}$, т.е.\
$\big(\frac{T_2}{T_1}\big)^\gamma=\frac{T_2}{T_3}$. \hfill$\blacksquare$
\end{problem}

%-------------------------------------------------------------
\begin{problem}{2.20}{dz}{Обратимый цикл: изотермическое расширение, изобарическое сжатие, адиабатическое сжатие (рис.~2.8).
Определить КПД, если отношение максимальной и минимальной абсолютных температур равно $\alpha$. $C_V$, $C_p$ постоянны.}
\textbf{1. Температуры.} Изотерма 1--2 — при $T_1=T_{\max}$. При изобарном сжатии 2--3 газ охлаждается до $T_3=T_{\min}=T_1/\alpha$; адиабата 3--1 возвращает к $T_1$.

\textbf{2. Энтропия} (за цикл $\Delta S=0$, на адиабате $0$): $\Delta S_{12}=-\Delta S_{23}=-C_p\ln\frac{T_3}{T_1}=C_p\ln\alpha$.

\textbf{3. Теплоты.} Изотерма: $Q_{12}=T_1\Delta S_{12}=C_pT_1\ln\alpha>0$ (получено). Изобара: $Q_{23}=C_p(T_3-T_1)<0$ (отдано).
\[\eta=1-\frac{|Q_{23}|}{Q_{12}}=1-\frac{C_p(T_1-T_1/\alpha)}{C_pT_1\ln\alpha}=1-\frac{1-1/\alpha}{\ln\alpha}.\]
\answer{$\eta=1-\dfrac{\alpha-1}{\alpha\ln\alpha}$.}
\end{problem}

%-------------------------------------------------------------
\begin{problem}{2.26}{dz}{Теплоизолированный цилиндр разделён пополам пористой перегородкой. Слева один моль водорода при $T_1=270$\,К, справа — моль азота при $T_2=290$\,К.
Перегородка пропускает водород и не пропускает азот. Найти $\Delta S$ после установления равновесия. Теплоёмкостью цилиндра пренебречь.}
\textbf{1. Конечное состояние.} Водород заполняет весь цилиндр ($V\to2V$), азот остаётся в своей половине.
Оба газа двухатомные, $C_V=\frac52R$; из сохранения энергии $\frac52R(T-T_1)+\frac52R(T-T_2)=0\Rightarrow T=\frac{T_1+T_2}2=280$\,К.

\textbf{2. Энтропия каждого газа:}
H$_2$: $\frac52R\ln\frac T{T_1}+R\ln2$ (нагрев + удвоение объёма); \;N$_2$: $\frac52R\ln\frac T{T_2}$ (только охлаждение).

\textbf{3. Сумма и числа:} $\frac52R\ln\frac{T^2}{T_1T_2}=\frac52\cdot8{,}31\cdot\ln\frac{280^2}{270\cdot290}=0{,}03$; $R\ln2=5{,}76$.
\answer{$\Delta S=\frac52R\ln\dfrac{(T_1+T_2)^2}{4T_1T_2}+R\ln2\approx5{,}8$ Дж/К (в книге $5{,}77$).}
\end{problem}

%-------------------------------------------------------------
\begin{problem}{2.27}{dz}{Сосуд объёмом $2V$ с теплонепроницаемыми стенками разделён пополам теплонепроницаемым поршнем. В каждой части по молю идеального газа
($C_V$ постоянна), начальные температуры $T_1$ и $T_2$. Поршень отпускают, его колебания затухают из-за внутреннего трения в газе.
После остановки поршень делит объём в отношении $V_1/V_2=n$. Трение поршня о стенки не учитывать. Найти $\Delta S$ газа.}
\textbf{1. Конечные объёмы:} $V_1'+V_2'=2V$, $V_1'/V_2'=n$ $\Rightarrow V_1'=\frac{2nV}{1+n}$, $V_2'=\frac{2V}{1+n}$.

\textbf{2. Конечные температуры.} Вся система изолирована, работа поршня — внутренняя, поэтому полная энергия газа сохраняется:
$C_V(T_1'+T_2')=C_V(T_1+T_2)$. Поршень остановился — давления равны: $\frac{RT_1'}{V_1'}=\frac{RT_2'}{V_2'}\Rightarrow\frac{T_1'}{T_2'}=n$. Отсюда
\[T_1'=\frac{n(T_1+T_2)}{1+n},\qquad T_2'=\frac{T_1+T_2}{1+n}.\]

\textbf{3. Энтропия:} $\Delta S=\sum_i\Big[C_V\ln\frac{T_i'}{T_i}+R\ln\frac{V_i'}{V}\Big]=C_V\ln\frac{T_1'T_2'}{T_1T_2}+R\ln\frac{V_1'V_2'}{V^2}$.
\answer{$\Delta S=C_V\ln\dfrac{n(T_1+T_2)^2}{(1+n)^2T_1T_2}+R\ln\dfrac{4n}{(1+n)^2}$.}
\end{problem}

%-------------------------------------------------------------
\begin{problem}{2.28}{dz}{Моль одноатомного идеального газа находится под поршнем в адиабатически изолированном цилиндре. Масса груза на поршне,
определяющая давление газа, мгновенно увеличивается вдвое. Найти изменение энтропии газа после установления равновесия. Атмосферным давлением пренебречь.}
\textbf{1. Процесс необратимый:} газ сжимается под постоянным внешним давлением $2p_0$ (груз). Работа над газом $2p_0(V_0-V)$, $Q=0$:
\[\tfrac32R(T-T_0)=2p_0(V_0-V).\]

\textbf{2. Конечное состояние.} Подставляем $p_0V_0=RT_0$ и $2p_0V=RT$ (в конце давление газа $2p_0$):
$\frac32R(T-T_0)=2RT_0-RT\Rightarrow T=\frac75T_0$. Объём: $\frac V{V_0}=\frac{T/2p_0}{T_0/p_0}=\frac{T}{2T_0}=0{,}7$.

\textbf{3. Энтропия} (формула через $T$ и $p$): $\Delta S=C_p\ln\frac T{T_0}-R\ln\frac p{p_0}=\frac52R\ln1{,}4-R\ln2=8{,}31\cdot(0{,}841-0{,}693)$.
\answer{$\Delta S=\frac52R\ln1{,}4-R\ln2\approx1{,}23$ Дж/К.}
\note{в книге то же в виде $\frac32R\ln2+\frac52R\ln0{,}7$.}
\end{problem}

%-------------------------------------------------------------
\begin{problem}{2.31}{dz}{В двух теплоизолированных сосудах находится по одному молю разных идеальных одноатомных газов при одинаковом давлении $p$.
Температура газа в первом сосуде $T_1$, во втором — $T_2$. Насколько изменится энтропия системы, если сосуды соединить?}
\textbf{1. Конечное состояние.} Объёмы $V_i=\frac{RT_i}{p}$. Энергия сохраняется: $\frac32R(T-T_1)+\frac32R(T-T_2)=0\Rightarrow T=\frac{T_1+T_2}2$.
Каждый газ занимает весь объём $V_1+V_2=\frac{R(T_1+T_2)}{p}$.

\textbf{2. Энтропия каждого газа:} $\Delta S_i=\frac32R\ln\frac{T}{T_i}+R\ln\frac{V_1+V_2}{V_i}=\frac32R\ln\frac{T}{T_i}+R\ln\frac{T_1+T_2}{T_i}$.

\textbf{3. Складываем:} $\Delta S=\frac32R\ln\frac{(T_1+T_2)^2}{4T_1T_2}+R\ln\frac{(T_1+T_2)^2}{T_1T_2}$, а
$R\ln\frac{(T_1+T_2)^2}{T_1T_2}=R\ln\frac{(T_1+T_2)^2}{4T_1T_2}+2R\ln2$.
\answer{$\Delta S=\frac52R\ln\dfrac{(T_1+T_2)^2}{4T_1T_2}+2R\ln2$.}
\end{problem}
